\title{طراحی آشکارساز شیء سبک و دقیق مبتنی بر YOLO با توجه به Transformer-in-Transformer و بلوک‌های Inception}
\name{محمدرضا}
\surname{سعادتی}
\degree{کارشناسی ارشد}
\thesisdate{۱۴۰۵شهریور }
\firstsupervisor{دکتر منصور رزقی}
\keywords{آشکارسازی شیء، YOLOv13، Transformer-in-Transformer، Inception، یادگیری عمیق}
\fa-abstract{این پایان‌نامه به طراحی یک آشکارساز شیء بلادرنگ می‌پردازد که هدف آن دستیابی هم‌زمان به دقت مناسب و هزینه محاسباتی پایین است. مبنای روش، YOLOv13 است که برای پردازش کارا از کانولوشن‌های تفکیک‌پذیر، تقویت هم‌بستگی‌های چندمقیاسی و توزیع ویژگی در سراسر خط لوله استفاده می‌کند. در معماری مورد مطالعه، توجه Transformer-in-Transformer برای مدل‌سازی وابستگی‌های درون‌قطعه‌ای و میان‌قطعه‌ای و بلوک‌های Inception همراه با توجه کانالی کارا برای استخراج ویژگی‌های چندمقیاسی به کار گرفته می‌شوند. همچنین گونه‌های توجه هندسی مبتنی بر نگاشت استریوگرافیک برای تحلیل تابع شباهت در بلوک‌های توجه بررسی می‌شوند. ارزیابی روش با معیارهای AP، AP50، AP75، تعداد پارامترها، FLOPs، تأخیر استنتاج و مصرف حافظه انجام می‌شود. سهم اصلی پژوهش، ارائه و ارزیابی یک طراحی سبک‌وزن است که ظرفیت بازنمایی ویژگی‌های محلی و چندمقیاسی را در آشکارسازی شیء تقویت می‌کند.}
\vtitle
